\chapter{Measures on Metric Spaces}

\begin{center}\itshape
How long is the coast of Britain? Measure it with a ruler a kilometre long, then with a ruler a metre long, then with a
centimetre: the answer keeps growing. Is the coastline a curve at all, or something between a curve and a surface?
\end{center}

The English meteorologist Lewis Fry Richardson, who spent his life trying to predict wars and the weather with
mathematics, left behind at his death in 1953 a strange table, published in 1961. Measuring the borders and coastlines of several countries on maps of increasing
precision, he had found that the measured length $L$ grows like a power of the inverse ruler size, $L\approx C\varepsilon^{1-D}$, with
$D\approx1.25$ for the west coast of Britain. Nobody paid attention until Beno\^it Mandelbrot published ``How long is the coast of
Britain?'' in \emph{Science} in 1967 and gave the exponent its meaning: $D$ is a \emph{dimension}, and coastlines have dimension
strictly between $1$ and $2$. The right notion had been invented half a century before, by Felix Hausdorff in 1918, with a measure
that can be computed in any metric space, in any real dimension.

This chapter is about measures that live on metric spaces, where distances give them extra structure. Three stories meet
here. The first is Frigyes Riesz's discovery (1909) that measures and integrals are two faces of one thing: a positive linear way of
averaging continuous functions \emph{is} a measure. The second is weak convergence, the convergence of laws that the central
limit theorem is about, with its portmanteau of equivalent forms and Prokhorov's compactness theorem (1956). The third is
Kolmogorov's extension theorem (1933), the last step in founding probability: a random sequence is determined by its finite
pieces, and any consistent choice of finite pieces is realised. Volume~4 admitted it; here it is proved. And at the end, Hausdorff
measures compute the dimension of the Cantor set, $\frac{\log2}{\log3}$, and of the carpet on the cover of this book.

\section{Regularity and tightness}

Throughout, $X$ is a metric space and its Borel sets are those of its topology. Lemma~\ref{lem:outerreg} showed that every finite Borel measure $\mu$ on $X$ is \emph{regular}: for every Borel
set $A$ and $\varepsilon>0$ there are a closed $C$ and an open $U$ with $C\subseteq A\subseteq U$ and $\mu(U\setminus C)<\varepsilon$. Compact sets are better than closed ones, and in the spaces probabilists use they
suffice. A metric space is \emph{Polish} if it is complete and separable: $\R^d$, $C([0,1])$, $\R^\N$ with its product metric, every closed subset of these.

\begin{theorem}[Ulam]\label{thm:ulam}
Every finite Borel measure $\mu$ on a Polish space $X$ is \emph{tight}: for every $\varepsilon>0$ there is a compact $K$ with $\mu(X\setminus K)<\varepsilon$. Hence $\mu(A)=\sup\{\mu(K):K\subseteq A\text{ compact}\}$ for every
Borel $A$.
\end{theorem}
\begin{proof}
Let $(x_i)$ be dense. For each $k\ge1$ the closed balls $\overline B(x_i,\frac1k)$ cover $X$, so by continuity from below there is $n_k$ with $\mu\bigl(X\setminus\bigcup_{i\le n_k}\overline B(x_i,\frac1k)\bigr)<\varepsilon2^{-k}$. The set
$K=\bigcap_k\bigcup_{i\le n_k}\overline B(x_i,\frac1k)$ is closed, hence complete, and totally bounded (covered by $n_k$ balls of radius $\frac1k$ for every $k$), hence compact (Theorem~\ref{thm:compactmetric}), and
$\mu(X\setminus K)<\varepsilon$. For the second statement, take $C\subseteq A$ closed with $\mu(A\setminus C)<\varepsilon$; then $C\cap K$ is compact.
\end{proof}

So a probability on a Polish space is, up to $\varepsilon$, carried by a compact set, even when the space is as large as $C([0,1])$: this is what makes Brownian motion (Volume~12) manageable.

\section{The Riesz representation theorem}

Every finite Borel measure $\mu$ on a compact metric space $K$ defines a linear form $\Lambda f=\int f\,d\mu$ on $C(K)$, which is \emph{positive}: $f\ge0$ implies $\Lambda f\ge0$. Riesz's theorem says that there are no other
positive linear forms. We need a partition of unity. The \emph{support} of $f$ is the closure of $\{f\ne0\}$, and $f\prec U$ means that $f\in C(K)$, $0\le f\le1$ and the support of $f$ lies in the open set $U$.

\begin{lemma}[partitions of unity]\label{lem:pou}
Let $C\subseteq K$ be closed and $C\subseteq V_1\cup\cdots\cup V_n$ with $V_i$ open. There are $h_i\prec V_i$ with $h_1+\cdots+h_n=1$ on $C$.
\end{lemma}
\begin{proof}
Each $x\in C$ has a radius $r_x>0$ with $\overline B(x,2r_x)\subseteq V_i$ for some $i=i(x)$. Finitely many balls $B(x,r_x)$ cover the compact $C$; let $W_i$ be the union of those with $i(x)=i$. Then $\overline{W_i}$ is compact and inside
$V_i$, so $\delta_i=d(\overline{W_i},K\setminus V_i)>0$ (or $V_i=K$), and $U_i=\{y:d(y,\overline{W_i})<\frac{\delta_i}2\}$ has $\overline{U_i}\subseteq V_i$. The Urysohn function $g_i(y)=\frac{d(y,K\setminus U_i)}{d(y,K\setminus U_i)+d(y,\overline{W_i})}$
(Proposition~\ref{prop:urysohn}) equals $1$ on $\overline{W_i}$ and $g_i\prec V_i$. Put $h_1=g_1$, $h_2=(1-g_1)g_2$, \ldots, $h_n=(1-g_1)\cdots(1-g_{n-1})g_n$. Then $h_i\prec V_i$, and by induction
$h_1+\cdots+h_n=1-(1-g_1)\cdots(1-g_n)$, which is $1$ on $C\subseteq\bigcup_i\overline{W_i}$.
\end{proof}

\begin{theorem}[Riesz representation]\label{thm:riesz}
Let $K$ be a compact metric space and $\Lambda$ a positive linear form on $C(K,\R)$. There is a unique finite Borel measure $\mu$ on $K$ with $\Lambda f=\int f\,d\mu$ for all $f\in C(K)$.
\end{theorem}
\begin{proof}
\emph{Construction.} For $U$ open let $\mu(U)=\sup\{\Lambda f:f\prec U\}$, and for any $A\subseteq K$ let $\mu^*(A)=\inf\{\mu(U):U\supseteq A\text{ open}\}$. Positivity makes $\mu$ monotone on open sets, so $\mu^*=\mu$ on them.

\emph{$\mu^*$ is an outer measure.} It suffices to show $\mu(\bigcup_iU_i)\le\sum_i\mu(U_i)$ for open sets (then argue as in Theorem~\ref{thm:extension}). Let $f\prec U=\bigcup_iU_i$. Its support is compact, so it lies in
$U_1\cup\cdots\cup U_n$ for some $n$; Lemma~\ref{lem:pou} gives $h_i\prec U_i$ with $\sum_{i\le n}h_i=1$ on the support, so $f=\sum_ih_if$ with $h_if\prec U_i$, and $\Lambda f=\sum_i\Lambda(h_if)\le\sum_i\mu(U_i)$.

\emph{Open sets are measurable.} By Carath\'eodory's criterion we need $\mu^*(A)\ge\mu^*(A\cap U)+\mu^*(A\setminus U)$, and by the definition of $\mu^*$ it suffices to take $A=V$ open. Given $\varepsilon$, pick $f\prec V\cap U$ with
$\Lambda f>\mu(V\cap U)-\varepsilon$, then $g\prec V\setminus\operatorname{supp}f$ (an open set containing $V\setminus U$) with $\Lambda g>\mu(V\setminus\operatorname{supp}f)-\varepsilon\ge\mu^*(V\setminus U)-\varepsilon$. Since $f+g\prec V$,
$\mu(V)\ge\Lambda f+\Lambda g\ge\mu(V\cap U)+\mu^*(V\setminus U)-2\varepsilon$. So $\mu^*$ restricts to a Borel measure $\mu$ (Theorem~\ref{thm:caratheodory}), with $\mu(K)=\Lambda1$ (as $f\prec K$ means $0\le f\le1$).

\emph{$\Lambda f=\int f\,d\mu$.} It suffices to prove $\Lambda f\le\int f\,d\mu$ for every real $f$, then apply it to $-f$. Let $f(K)\subseteq\,]a,b[$, $\varepsilon>0$, and $y_0<a<y_1<\cdots<y_n=b$ with $y_i-y_{i-1}<\varepsilon$. The sets
$E_i=\{y_{i-1}<f\le y_i\}$ are disjoint Borel sets covering $K$. Choose open $V_i\supseteq E_i$ with $\mu(V_i)<\mu(E_i)+\frac\varepsilon n$ and $f<y_i+\varepsilon$ on $V_i$ (intersect with $\{f<y_i+\varepsilon\}$), and $h_i\prec V_i$ with $\sum_ih_i=1$
on $K$ (Lemma~\ref{lem:pou}). Then $\Lambda h_i\le\mu(V_i)$, $\sum_i\Lambda h_i=\Lambda1=\mu(K)$, and $h_if\le(y_i+\varepsilon)h_i$, so with $c=|a|$,
\[\begin{aligned}\Lambda f=\sum_i\Lambda(h_if)&\le\sum_i(c+y_i+\varepsilon)\Lambda h_i-c\,\mu(K)\eqbr\le\sum_i(c+y_i+\varepsilon)\Bigl(\mu(E_i)+\frac\varepsilon n\Bigr)-c\sum_i\mu(E_i)\\&\le\sum_i(y_i-\varepsilon)\mu(E_i)+2\varepsilon\mu(K)+\varepsilon(c+b+\varepsilon)\eqbr\le\int f\,d\mu+\varepsilon\bigl(2\mu(K)+c+b+\varepsilon\bigr),\end{aligned}\]
using $c+y_i+\varepsilon\ge0$ and $y_i-\varepsilon\le f$ on $E_i$. Let $\varepsilon\to0$.

\emph{Uniqueness.} If $\int f\,d\mu=\int f\,d\nu$ for all $f\in C(K)$, apply it to $f_n(x)=\max(0,1-n\,d(x,C))$, which decreases to $\mathbf 1_C$ for $C$ closed: $\mu(C)=\nu(C)$ by dominated convergence. Closed sets form a
$\pi$-system generating the Borel sets, and $\mu(K)=\nu(K)$, so $\mu=\nu$ (Theorem~\ref{thm:uniq6}).
\end{proof}

Riesz proved this in 1909 for $K=[0,1]$, where it says that every positive linear form on $C([0,1])$ is a Stieltjes integral $\int f\,dF$ with $F$ non-decreasing. The theorem turns measure theory upside down:
one may start from the integral, a linear form on continuous functions, and recover the measure. Volume~8 extends it to signed measures, which makes the space of finite measures the dual of $C(K)$.

\section{Weak convergence}

A sequence of probabilities $\mu_n$ on $X$ \emph{converges weakly} to $\mu$, written $\mu_n\Rightarrow\mu$, if $\int f\,d\mu_n\to\int f\,d\mu$ for every bounded continuous $f$. For random variables, $X_n\Rightarrow X$ means that
their laws converge weakly; this is convergence in distribution, as in the central limit theorem of Volume~4.

\begin{think}
If $\mu_n\Rightarrow\mu$, must $\mu_n(A)\to\mu(A)$ for every Borel set $A$? Try $\mu_n=\delta_{1/n}$ and $A=\{0\}$, or $A=]0,1]$.
\end{think}

$\delta_{1/n}\Rightarrow\delta_0$, since $f(\frac1n)\to f(0)$ for continuous $f$; but $\delta_{1/n}(\{0\})=0\not\to1$ and $\delta_{1/n}(]0,1])=1\not\to0$. Mass can move onto the boundary of a set at the last moment, and the boundary is
exactly where things go wrong. The \emph{boundary} of $A$ is $\partial A=\overline A\setminus\mathring A$, and $A$ is a \emph{$\mu$-continuity set} if $\mu(\partial A)=0$.

\begin{theorem}[portmanteau]\label{thm:portmanteau}
For probabilities $\mu_n,\mu$ on a metric space $X$, the following are equivalent:
\begin{enumerate}[label=(\alph*)]
\item $\mu_n\Rightarrow\mu$;
\item $\int f\,d\mu_n\to\int f\,d\mu$ for every bounded Lipschitz $f$;
\item $\limsup_n\mu_n(F)\le\mu(F)$ for every closed $F$;
\item $\liminf_n\mu_n(U)\ge\mu(U)$ for every open $U$;
\item $\mu_n(A)\to\mu(A)$ for every Borel $\mu$-continuity set $A$.
\end{enumerate}
\end{theorem}
\begin{proof}
(a)$\Rightarrow$(b) is clear. (b)$\Rightarrow$(c): $f_k(x)=\max(0,1-k\,d(x,F))$ is Lipschitz, $\mathbf 1_F\le f_k$, and $f_k\downarrow\mathbf 1_F$. So $\limsup_n\mu_n(F)\le\lim_n\int f_k\,d\mu_n=\int f_k\,d\mu$, which tends to $\mu(F)$ as
$k\to\infty$. (c)$\Leftrightarrow$(d) by complements. (c),(d)$\Rightarrow$(e): $\mu(\mathring A)\le\liminf\mu_n(\mathring A)\le\liminf\mu_n(A)\le\limsup\mu_n(A)\le\limsup\mu_n(\overline A)\le\mu(\overline A)$, and the two ends are equal to
$\mu(A)$ when $\mu(\partial A)=0$. (e)$\Rightarrow$(a): let $f$ be continuous with $0<f<1$ (after an affine change). For all but countably many $t$, $\mu(\{f=t\})=0$ (these sets are disjoint), and $\partial\{f>t\}\subseteq\{f=t\}$ by continuity.
So $\mu_n(f>t)\to\mu(f>t)$ for almost every $t$, and by the layer-cake formula (Exercise~\ref{ex:layercake}) and dominated convergence on $[0,1]$,
$\int f\,d\mu_n=\int_0^1\mu_n(f>t)\,dt\to\int_0^1\mu(f>t)\,dt=\int f\,d\mu$.
\end{proof}

\begin{corollary}\label{cor:cdfweak}
Probabilities on $\R$ converge weakly iff their distribution functions converge at every continuity point of the limit: $F_n(t)\to F(t)$ whenever $F$ is continuous at $t$.
\end{corollary}
\begin{proof}
$\partial\,]-\infty,t]=\{t\}$, so (e) gives the convergence at continuity points. Conversely, if $F_n\to F$ on the dense set $D$ of continuity points, every open set is a countable disjoint union of open
intervals, and each $]a,b[$ contains $]a',b']$ with $a',b'\in D$ and $\mu(]a,b[\setminus\,]a',b'])$ small; then $\liminf\mu_n(]a,b[)\ge\lim(F_n(b')-F_n(a'))=F(b')-F(a')$, close to $\mu(]a,b[)$. Fatou's lemma for sums
extends this to countable unions, giving (d).
\end{proof}

\begin{figure}[ht]
\centering
\fig{ch11-empirical}
\caption{Empirical distribution functions of $n=10$, $100$ and $1000$ uniform samples on $[0,1]$, converging to $F(t)=t$. The empirical laws $\frac1n\sum_i\delta_{U_i}$ converge weakly to the uniform law with probability
$1$ (Exercise~\ref{ex:glivenko}); each of them is carried by a finite, hence null, set.}
\label{fig:empirical}
\end{figure}

\section{Prokhorov's theorem}

When does a sequence of laws have a weakly convergent subsequence? Mass must not escape to infinity: a family $\mathcal M$ of probabilities is \emph{tight} if for every $\varepsilon>0$ there is a compact $K$ with $\mu(K)\ge1-\varepsilon$ for all
$\mu\in\mathcal M$. The laws $\delta_n$ on $\R$ are not tight, and indeed have no weakly convergent subsequence.

\begin{lemma}\label{lem:separable}
If $K$ is a compact metric space, $C(K)$ is separable.
\end{lemma}
\begin{proof}
Let $(a_k)$ be dense in $K$ (Chapter~4). The algebra generated over $\Q$ by $1$ and the functions $d(\cdot,a_k)$ is countable. Its closure is a closed subalgebra (closed under real multiples, as limits of rational ones),
contains the constants, and separates points: if $x\ne y$, pick $a_k$ with $d(x,a_k)<\frac13d(x,y)$, so $d(x,a_k)\ne d(y,a_k)$. By Stone--Weierstrass (Theorem~\ref{thm:sw}) it is $C(K)$.
\end{proof}

\begin{theorem}[Prokhorov]\label{thm:prokhorov}
A tight sequence of probabilities on a compact metric space, or on $\R^d$, has a weakly convergent subsequence.
\end{theorem}
\begin{proof}
\emph{Compact $K$.} Every sequence is tight. Let $(f_j)$ be dense in $C(K)$. Each $(\int f_j\,d\mu_n)_n$ is bounded by $\|f_j\|_\infty$, so a diagonal extraction gives a subsequence, still written $\mu_n$, along which $\int f_j\,d\mu_n$
converges for every $j$. For $f\in C(K)$ and $\|f-f_j\|_\infty<\varepsilon$, the sequence $\int f\,d\mu_n$ is within $2\varepsilon$ of a convergent one, so it is Cauchy: $\Lambda f=\lim\int f\,d\mu_n$ exists. $\Lambda$ is linear, positive and
$\Lambda1=1$, so by Riesz $\Lambda f=\int f\,d\mu$ for a probability $\mu$, and $\mu_n\Rightarrow\mu$.

\emph{$\R^d$.} The map $\varphi(x)=(\frac{x_i}{1+|x_i|})_i$ is a homeomorphism of $\R^d$ onto the open cube $Q=\,]-1,1[^d$, inside the compact $\overline Q$. Let $\nu_n=\varphi_*\mu_n$, probabilities on $\overline Q$; a subsequence converges
weakly to some $\nu$ on $\overline Q$. For $\varepsilon>0$ let $K$ be compact with $\mu_n(K)\ge1-\varepsilon$ for all $n$; $\varphi(K)$ is compact, and by the portmanteau theorem $\nu(\varphi(K))\ge\limsup\nu_n(\varphi(K))\ge1-\varepsilon$. So
$\nu(Q)=1$, and $\mu=\varphi^{-1}_*\nu$ is a probability on $\R^d$. For $g\in C_c(\R^d)$, $g\circ\varphi^{-1}$ extends by $0$ to a continuous function on $\overline Q$, so $\int g\,d\mu_n\to\int g\,d\mu$ along the subsequence. To pass to
bounded continuous $f$, take $\chi\in C_c$ with $0\le\chi\le1$ and $\int\chi\,d\mu>1-\varepsilon$; then $\int(1-\chi)\,d\mu_n<\varepsilon$ for $n$ large, and
\[\Bigl|\int f\,d\mu_n-\int f\,d\mu\Bigr|\le\Bigl|\int f\chi\,d\mu_n-\int f\chi\,d\mu\Bigr|+\|f\|_\infty\Bigl(\int(1-\chi)\,d\mu_n+\int(1-\chi)\,d\mu\Bigr),\]
which is eventually less than $(2\|f\|_\infty+1)\varepsilon$, since $f\chi\in C_c$.
\end{proof}

Prokhorov proved more (1956): on a Polish space, a family is tight if and only if every sequence in it has a weakly convergent subsequence. This is the compactness behind Donsker's theorem, the
central limit theorem for random walks seen as random functions (Volume~12), and Exercise~\ref{ex:helly} gives Helly's original one-dimensional version.

\section{Kolmogorov's extension theorem}

A random sequence $(X_1,X_2,\ldots)$ is a random point of $\R^\N$. Its finite-dimensional laws $\mu_n=$ law of $(X_1,\dots,X_n)$ are \emph{consistent}: $\mu_{n+1}(B\times\R)=\mu_n(B)$ for every Borel $B\subseteq\R^n$. Kolmogorov's
theorem says that consistency is the only condition. On $\R^\N$, let $\pi_n(x)=(x_1,\dots,x_n)$; a \emph{cylinder} is a set $\pi_n^{-1}(B)$ with $B$ Borel in $\R^n$, and the product $\sigma$-algebra is generated by the cylinders.

\begin{theorem}[Kolmogorov, 1933]\label{thm:kolmogorov}
Let $\mu_n$ be consistent probabilities on $\R^n$, $n\ge1$. There is a unique probability $\mu$ on the product $\sigma$-algebra of $\R^\N$ with $\mu(\pi_n^{-1}(B))=\mu_n(B)$ for all $n$ and Borel $B\subseteq\R^n$.
\end{theorem}
\begin{proof}
The cylinders form an algebra $\mathcal C$: two cylinders can be written with the same $n$, by consistency of the description $\pi_n^{-1}(B)=\pi_{n+1}^{-1}(B\times\R)$. The set function $\mu_0(\pi_n^{-1}(B))=\mu_n(B)$ is well defined by consistency: if $\pi_n^{-1}(B)=\pi_m^{-1}(B')$ with $n\le m$, then $\mu_m(B\times\R^{m-n})=\mu_n(B)$. It is finitely additive. A finite family of cylinders may be written at one common level $N$ as $\pi_N^{-1}(B_j)$; if the cylinders are disjoint then the $B_j$ are disjoint, and if their union is a cylinder $\pi_N^{-1}(B)$ then $B=\bigsqcup_j B_j$, so $\mu_0$ of the union is $\mu_N(B)=\sum_j\mu_N(B_j)$. An algebra is a semiring, so by Carath\'eodory's extension theorem (Theorem~\ref{thm:extension}) it suffices to show that $\mu_0$ is countably additive on $\mathcal C$; uniqueness then follows from Theorem~\ref{thm:uniq6}.
By finite additivity, it is enough to show: if cylinders $C_k\downarrow\varnothing$, then $\mu_0(C_k)\to0$.

Suppose instead that $\mu_0(C_k)\ge\delta>0$ for all $k$, with $C_k=\pi_{n_k}^{-1}(B_k)$, $n_1<n_2<\cdots$. Each $\mu_{n_k}$ is tight (Theorem~\ref{thm:ulam}), so there is a compact $K_k\subseteq B_k$ with
$\mu_{n_k}(B_k\setminus K_k)<\delta2^{-k-1}$. Let $D_k=\bigcap_{j\le k}\pi_{n_j}^{-1}(K_j)\subseteq C_k$. Then $\mu_0(C_k\setminus D_k)\le\sum_{j\le k}\mu_0(C_j\setminus\pi_{n_j}^{-1}(K_j))<\frac\delta2$, so $\mu_0(D_k)>\frac\delta2$ and $D_k\ne\varnothing$.
Pick $x^{(k)}\in D_k$. For each $j$, the points $\pi_{n_j}(x^{(k)})$, $k\ge j$, lie in the compact $K_j$; a diagonal extraction gives a subsequence along which every coordinate of $x^{(k)}$ converges, to some $x\in\R^\N$, and
$\pi_{n_j}(x)\in K_j$ for every $j$ since $K_j$ is closed. So $x\in\bigcap_j\pi_{n_j}^{-1}(K_j)\subseteq\bigcap_jC_j=\varnothing$: a contradiction.
\end{proof}

With $\mu_n=\nu_1\otimes\cdots\otimes\nu_n$ this gives independent sequences with any laws, a second proof of Theorem~\ref{thm:independent}. But the theorem builds much more: Markov chains from transition kernels, Gaussian
sequences from covariances, any process indexed by $\N$ whose finite-dimensional laws are known. Volume~4 admitted it, and Volume~12 extends it to processes indexed by $[0,\infty[$, where it is the first step in constructing
Brownian motion. The product $\sigma$-algebra is the Borel $\sigma$-algebra of $\R^\N$ with the product topology, a Polish space (as in Proposition~\ref{prop:borelproduct}).

\section{Hausdorff measures and dimension}

\begin{figure}[ht]
\centering
\fig{ch11-carpet}
\caption{The Sierpi\'nski carpet, the attractor of eight similarities of ratio $\frac13$ (Chapter~2). Each level of the construction keeps $8$ of the $9$ subsquares. Its Hausdorff dimension is $\frac{\log8}{\log3}\approx1.893$
(Exercise~\ref{ex:carpet}).}
\label{fig:carpet}
\end{figure}

Lengths, areas and volumes scale like $r^1$, $r^2$, $r^3$ under a dilation by $r$. Hausdorff's idea is to measure a set by covering it with small pieces and adding up $(\text{diameter})^s$, for any real $s\ge0$. For $A\subseteq X$
and $\delta>0$ let
\[\mathcal H^s_\delta(A)=\inf\Bigl\{\sum_i(\operatorname{diam}U_i)^s:A\subseteq\bigcup_iU_i,\ \operatorname{diam}U_i\le\delta\Bigr\},\]
over countable covers, with $(\operatorname{diam}\varnothing)^s=0$. As $\delta$ decreases, fewer covers are allowed and $\mathcal H^s_\delta(A)$ increases; its limit $\mathcal H^s(A)\in[0,\infty]$ is the \emph{$s$-dimensional Hausdorff measure}. Like
$\lambda^*$, $\mathcal H^s_\delta$ and $\mathcal H^s$ are outer measures (the proof of Theorem~\ref{thm:extension} applies word for word).

\begin{remark}
$\mathcal H^s$ restricted to the Borel sets is a measure, because it is a \emph{metric} outer measure ($\mathcal H^s(A\cup B)=\mathcal H^s(A)+\mathcal H^s(B)$ when $d(A,B)>0$), and on $\R^d$, $\mathcal H^d$ is a constant multiple of $\lambda_d$.
Proofs are in Falconer's book (Sources); we shall only need $\mathcal H^s$ as an outer measure.
\end{remark}

\begin{proposition}\label{prop:jump}
If $\mathcal H^s(A)<\infty$ and $t>s$, then $\mathcal H^t(A)=0$. Hence there is a unique $\dim_HA\in[0,\infty]$, the \emph{Hausdorff dimension}, with $\mathcal H^s(A)=\infty$ for $s<\dim_HA$ and $\mathcal H^s(A)=0$ for $s>\dim_HA$.
\end{proposition}
\begin{proof}
For a cover by sets of diameter at most $\delta$, \[\sum_i(\operatorname{diam}U_i)^t\le\delta^{t-s}\sum_i(\operatorname{diam}U_i)^s,\] so $\mathcal H^t_\delta(A)\le\delta^{t-s}\mathcal H^s_\delta(A)\le\delta^{t-s}\mathcal H^s(A)\to0$.
\end{proof}

\begin{think}
What is the dimension of a set that is ``between'' a line and a plane? Can a set of Lebesgue measure zero in the plane have dimension $2$?
\end{think}

Upper bounds on dimension come from explicit covers. Lower bounds are harder, since they concern \emph{all} covers, and the trick is to spread a measure on the set.

\begin{lemma}[mass distribution principle]\label{lem:massdist}
Let $\mu$ be a finite Borel measure with $\mu(A)>0$, and suppose $\mu(U)\le c\,(\operatorname{diam}U)^s$ for every closed set $U$ of diameter at most $\delta$. Then $\mathcal H^s(A)\ge\frac{\mu(A)}c>0$, so $\dim_HA\ge s$.
\end{lemma}
\begin{proof}
If $A\subseteq\bigcup_iU_i$ with $\operatorname{diam}U_i\le\delta$, replace each $U_i$ by its closure (same diameter); then \[\mu(A)\le\sum_i\mu(\overline{U_i})\le c\sum_i(\operatorname{diam}U_i)^s.\] Take the infimum over covers.
\end{proof}

\begin{figure}[ht]
\centering
\fig{ch11-cantorcover}
\caption{Covering the Cantor set at three scales: $2$ intervals of length $\frac13$, $4$ of length $\frac19$, $8$ of length $\frac1{27}$. With $s=\frac{\log2}{\log3}$, each cover has $\sum(\operatorname{diam})^s=2^n3^{-ns}=1$: the
number of pieces and their size balance exactly at this exponent.}
\label{fig:cantorcover}
\end{figure}

\begin{theorem}\label{thm:cantordim}
The Cantor set has Hausdorff dimension $s=\frac{\log2}{\log3}\approx0.631$, and $\frac12\le\mathcal H^s(C)\le1$.
\end{theorem}
\begin{proof}
\emph{Upper bound.} The $2^n$ intervals of $C_n$ have length $3^{-n}$ and cover $C$, and $2^n(3^{-n})^s=1$ since $3^s=2$. So $\mathcal H^s_{3^{-n}}(C)\le1$ for every $n$, and $\mathcal H^s(C)\le1$.

\emph{Lower bound.} Let $\mu_c$ be the Cantor law (Chapter~7), which gives mass $2^{-k}$ to each of the $2^k$ intervals of $C_k$ and lives on $C$. Two intervals of $C_k$ are separated by a gap of length at least $3^{-k}$.
Let $U$ be closed with $0<\operatorname{diam}U<1$, and $k$ such that $3^{-k-1}\le\operatorname{diam}U<3^{-k}$. Then $U$ meets at most one interval of $C_k$, so
$\mu_c(U)\le2^{-k}=(3^{-k})^s=3^s(3^{-k-1})^s\le2(\operatorname{diam}U)^s$. A set of diameter $0$ is a point, of measure $0$. Lemma~\ref{lem:massdist} with $c=2$ gives $\mathcal H^s(C)\ge\frac12$.
\end{proof}

So the Cantor set is too thin to have a length ($\mathcal H^1(C)=0$, as $1>s$) and too fat to be counted ($\mathcal H^0$ is counting measure, and $C$ is uncountable); in dimension $0.631$ its size is a positive, finite number.
In fact $\mathcal H^s(C)=1$, a theorem of Hausdorff that needs a sharper lower bound. The same argument gives the dimension of every set built by $m$ similar copies of ratio $r$ that do not overlap too much:
$\frac{\log m}{\log(1/r)}$ (Moran 1946, Hutchinson 1981). The exercises do the numbers without a digit $7$ and the Sierpi\'nski carpet.

\begin{beauty}
A measure is a way of saying how big a set is. Lebesgue measure answers for lines, planes and space; Hausdorff measures answer in every real dimension, and the answer for the Cantor set, the coast of
Britain or the carpet on the cover is a dimension that is not a whole number. On the probabilistic side, Riesz, Prokhorov and Kolmogorov turn measures into the points of a space with its own topology, where laws
converge, subsequences exist, and random sequences are built from their finite pieces. The stage is set for Volumes~11 and~12.
\end{beauty}

\section{Exercises}

\begin{exercise}[$\star$]\label{ex:nodigit7}
Let $E$ be the set of $x\in[0,1]$ having a decimal expansion without the digit $7$. Show that $E$ is closed, that $\lambda(E)=0$, and that $\dim_HE=\frac{\log9}{\log10}$. (Adjacent intervals of the construction may touch: count how many a
small set can meet.)
\end{exercise}

\begin{exercise}[$\star\star$]\label{ex:carpet}
Show that the Sierpi\'nski triangle has dimension $\frac{\log3}{\log2}$ and the Sierpi\'nski carpet $\frac{\log8}{\log3}$. Give an example of a set in the plane of dimension $2$ and Lebesgue measure $0$.
\end{exercise}

\begin{exercise}[$\star$]\label{ex:lipschitzdim}
Let $f:A\to Y$ be $L$-Lipschitz. Show that $\mathcal H^s(f(A))\le L^s\mathcal H^s(A)$ and $\dim_Hf(A)\le\dim_HA$. Deduce that the Cantor set and $[0,1]$ are not bi-Lipschitz equivalent, although there is a continuous surjection from
one onto the other (Exercise~\ref{ex:lebesguecurve}). Show that $\dim_H$ of a countable set is $0$ and that $\dim_H(\bigcup_kA_k)=\sup_k\dim_HA_k$.
\end{exercise}

\begin{exercise}[$\star\star$]\label{ex:helly}
\emph{Helly's selection theorem.} Let $F_n$ be distribution functions on $\R$. (a) Show that a subsequence converges at every rational to a limit $G$, and that $F(t)=\inf_{q>t}G(q)$ is non-decreasing and right-continuous
with $F_{n_k}(t)\to F(t)$ at every continuity point of $F$. (b) Show that $F$ is a distribution function (limits $0$ and $1$) if the laws are tight, and give an example where it is not.
\end{exercise}

\begin{exercise}[$\star\star$]\label{ex:glivenko}
\emph{Empirical laws.} Let $U_1,U_2,\ldots$ be independent uniform on $[0,1]$ and $F_n(t)=\frac1n\#\{i\le n:U_i\le t\}$. Using the strong law of large numbers at the points $\frac jm$ and monotonicity, show that
$\sup_t|F_n(t)-t|\to0$ almost surely (Glivenko--Cantelli). Deduce that $\frac1n\sum_{i\le n}\delta_{U_i}\Rightarrow\lambda|_{[0,1]}$ almost surely.
\end{exercise}

\begin{exercise}[$\star\star$]\label{ex:weakexamples}
(a) Show that $\mathcal N(0,\frac1n)\Rightarrow\delta_0$, and that the uniform laws on $\{\frac1n,\frac2n,\dots,1\}$ converge weakly to $\lambda|_{[0,1]}$. (b) Find laws $\mu_n\Rightarrow\mu$ on $\R$ with $\int x\,d\mu_n\not\to\int x\,d\mu$. (c) Show that
if $\mu_n\Rightarrow\mu$ and $\sup_n\int|x|^2\,d\mu_n<\infty$, then $\int x\,d\mu_n\to\int x\,d\mu$ (compare Vitali's theorem).
\end{exercise}

\section{Hints}

\begin{hint}{ex:nodigit7}
At stage $k$ there are $9^k$ closed intervals of length $10^{-k}$. A closed set of diameter less than $10^{-k}$ meets at most two of them. The law of the random decimal with independent digits uniform on
$\{0,\ldots,9\}\setminus\{7\}$ spreads mass $9^{-k}$ on each.
\end{hint}
\begin{hint}{ex:carpet}
At stage $k$ the carpet is covered by $8^k$ squares of side $3^{-k}$; a set of diameter less than $3^{-k}$ meets at most $4$ of them. For the last question, take a countable union of compact sets of dimensions
tending to $2$.
\end{hint}
\begin{hint}{ex:lipschitzdim}
Map a cover of $A$ to a cover of $f(A)$. For countable unions, $\mathcal H^s$ is countably subadditive.
\end{hint}
\begin{hint}{ex:helly}
(a) Diagonal extraction over the rationals. If $F$ is continuous at $t$, compare $F_{n_k}(t)$ with $F_{n_k}(q)$ for rationals $q$ close to $t$ on each side. (b) $F_n=\mathbf 1_{[n,\infty[}$.
\end{hint}
\begin{hint}{ex:glivenko}
For $t\in[\frac{j-1}m,\frac jm]$, $F_n(\frac{j-1}m)-\frac jm\le F_n(t)-t\le F_n(\frac jm)-\frac{j-1}m$.
\end{hint}
\begin{hint}{ex:weakexamples}
(a) Use Corollary~\ref{cor:cdfweak}, or Riemann sums. (b) Put mass $\frac1n$ at $n$. (c) Truncate: $|\int x\,d\mu_n-\int x\wedge M\vee(-M)\,d\mu_n|\le\frac1M\int x^2d\mu_n$.
\end{hint}

\section{Solutions}

\begin{solution}{ex:nodigit7}
\idea Same structure as the Cantor set, with $9$ pieces of ratio $\frac1{10}$ instead of $2$ of ratio $\frac13$. The only new feature is that kept intervals can touch, so a small set may meet two of them, which costs a factor $2$ in
the mass distribution principle and nothing in the dimension.
\writeup Let $E_k$ be the union of the $9^k$ closed intervals $[\frac m{10^k},\frac{m+1}{10^k}]$ where $m$ has $k$ decimal digits none of which is $7$. Then $E=\bigcap_kE_k$: a number with an expansion avoiding $7$ is in every $E_k$;
conversely a point of every $E_k$ has, at each level, a chosen digit different from $7$, and these choices are consistent (nested intervals), giving an expansion avoiding $7$. So $E$ is closed, and $\lambda(E)\le\lambda(E_k)=(\frac9{10})^k\to0$.
Upper bound: $\mathcal H^s_{10^{-k}}(E)\le9^k10^{-ks}=1$ for $s=\frac{\log9}{\log10}$. Lower bound: let $\mu$ be the law of $\sum_jD_j10^{-j}$ with independent $D_j$ uniform on the nine digits other than $7$; it lives on $E$ and gives
mass $9^{-k}$ to each interval of $E_k$. If $U$ is closed with $10^{-k-1}\le\operatorname{diam}U<10^{-k}$, it meets at most two intervals of $E_k$ (they have length $10^{-k}$ and disjoint interiors), so
$\mu(U)\le2\cdot9^{-k}=2(10^{-k})^s\le2\cdot10^s(\operatorname{diam}U)^s=18(\operatorname{diam}U)^s$. By Lemma~\ref{lem:massdist}, $\mathcal H^s(E)\ge\frac1{18}$, so $\dim_HE=\frac{\log9}{\log10}\approx0.954$. Almost every number has a $7$ in its
expansion, yet those without one form a set of dimension close to $1$.
\end{solution}

\begin{solution}{ex:carpet}
\idea Count pieces and measure their size at each stage for the upper bound; spread the natural uniform mass for the lower bound, controlling how many pieces of the right size a small set can touch.
\writeup \emph{Carpet.} Stage $k$ consists of $8^k$ closed squares of side $3^{-k}$, diameter $\sqrt2\,3^{-k}$, covering the carpet $S$; with $s=\frac{\log8}{\log3}$, $\mathcal H^s_{\sqrt23^{-k}}(S)\le8^k(\sqrt2\,3^{-k})^s=2^{s/2}$, so $\mathcal H^s(S)<\infty$.
Let $\mu$ be the law of the random point whose base-$3$ digit pairs are chosen independently and uniformly among the $8$ allowed pairs; it gives mass $8^{-k}$ to each square of stage $k$. A closed $U$ with
$3^{-k-1}\le\operatorname{diam}U<3^{-k}$ lies in a square of side $3^{-k}$ placed anywhere, which meets at most $4$ squares of the stage-$k$ grid. So $\mu(U)\le4\cdot8^{-k}=4(3^{-k})^s\le4\cdot3^s(\operatorname{diam}U)^s$, and
$\mathcal H^s(S)>0$. \emph{Triangle.} Stage $k$ has $3^k$ triangles of side $2^{-k}$, and the same argument with $s=\frac{\log3}{\log2}$ (a set of diameter less than $2^{-k}$ meets a bounded number of them, at most $6$ say,
since they lie in the triangular grid of side $2^{-k}$) gives $0<\mathcal H^s<\infty$. \emph{Dimension $2$, measure $0$.} Let $S_m$ be the generalised carpet removing $1$ central square out of $m^2$, at every stage, so
$\dim S_m=\frac{\log(m^2-1)}{\log m}\to2$, and each $S_m$ is null (its area is multiplied by $1-m^{-2}$ at each stage). Then $\bigcup_mS_m$ is null, and has dimension $\sup_m\dim S_m=2$ (Exercise~\ref{ex:lipschitzdim}).
\end{solution}

\begin{solution}{ex:lipschitzdim}
\idea A Lipschitz map shrinks diameters by at most $L$, so it maps covers to covers with sums multiplied by at most $L^s$. Dimension cannot go up.
\writeup If $A\subseteq\bigcup_iU_i$ with $\operatorname{diam}U_i\le\delta$, then $f(A)\subseteq\bigcup_if(U_i\cap A)$, with $\operatorname{diam}f(U_i\cap A)\le L\operatorname{diam}U_i\le L\delta$, so $\mathcal H^s_{L\delta}(f(A))\le L^s\sum_i(\operatorname{diam}U_i)^s$; take the
infimum and let $\delta\to0$. If $s>\dim_HA$, $\mathcal H^s(A)=0$, so $\mathcal H^s(f(A))=0$ and $\dim_Hf(A)\le s$. A bi-Lipschitz bijection would preserve dimension, but $\dim_HC=0.631\ne1=\dim_H[0,1]$ ($0<\mathcal H^1([0,1])<\infty$: it is
at most $1$ by covering with small intervals, and at least $\lambda([0,1])=1$ because $\lambda^*\le\mathcal H^1_\delta$, an interval of diameter $\ell$ having length $\ell$). A point has $\mathcal H^s=0$ for $s>0$, and $\mathcal H^s$ is countably
subadditive, so a countable set has $\mathcal H^s=0$ for all $s>0$. For a countable union: monotonicity gives $\ge$; if $s>\sup_k\dim_HA_k$, all $\mathcal H^s(A_k)=0$, so $\mathcal H^s(\bigcup A_k)=0$.
\end{solution}

\begin{solution}{ex:helly}
\idea Distribution functions are monotone and bounded, so countably many values can be made to converge by a diagonal argument, and monotonicity fills in the rest, except at jumps of the limit. Mass can
escape to infinity, and tightness is exactly what forbids it.
\writeup (a) The values $F_n(q)\in[0,1]$, $q\in\Q$, can be made to converge along one subsequence by diagonal extraction; let $G(q)$ be the limits, non-decreasing in $q$. $F(t)=\inf_{q>t}G(q)$ is non-decreasing and right-continuous
(an infimum over $q>t$ is the limit of the infima over $q>t'$ as $t'\downarrow t$). Let $F$ be continuous at $t$ and $\varepsilon>0$. By the definition of $F(t)$ there is a rational $q>t$ with $G(q)<F(t)+\varepsilon$. By continuity there is $s<t$ with $F(s)>F(t)-\varepsilon$; take a rational $r\in\,]s,t[$, so that $G(r)\ge\inf_{q'>s}G(q')=F(s)$. Monotonicity gives $F_{n_k}(r)\le F_{n_k}(t)\le F_{n_k}(q)$, and letting $k\to\infty$, $F(t)-\varepsilon<G(r)\le\liminf_kF_{n_k}(t)\le\limsup_kF_{n_k}(t)\le G(q)<F(t)+\varepsilon$. (b) If $\mu_n([-M,M])\ge1-\varepsilon$ for all $n$, then $F_n(-M-1)\le\varepsilon$ and
$F_n(M)\ge1-\varepsilon$, so the limits satisfy $F(t)\le\varepsilon$ for $t<-M-1$ and $F(t)\ge1-\varepsilon$ for $t>M$ (at continuity points, dense): $F$ goes from $0$ to $1$. With $F_n=\mathbf 1_{[n,\infty[}$ (the law $\delta_n$), $F=0$: all the
mass has escaped to $+\infty$.
\end{solution}

\begin{solution}{ex:glivenko}
\idea The strong law handles one point at a time; countably many points are no harder, since a countable union of null sets is null. Monotonicity of $F_n$ and of $t$ then controls all $t$ at once, up to the grid
size $\frac1m$.
\writeup For each rational $t$, $F_n(t)=\frac1n\sum_{i\le n}\mathbf 1_{U_i\le t}\to t$ almost surely by the strong law (Volume~4). Outside a null set this holds for all rational $t$ simultaneously. Fix $m$ and $t\in[\frac{j-1}m,\frac jm]$; by
monotonicity, $F_n(t)-t\le F_n(\frac jm)-\frac jm+\frac1m$ and $F_n(t)-t\ge F_n(\frac{j-1}m)-\frac{j-1}m-\frac1m$. So $\sup_t|F_n(t)-t|\le\max_{j\le m}|F_n(\frac jm)-\frac jm|+\frac1m$, whose $\limsup$ is at most $\frac1m$; let $m\to\infty$. Weak
convergence follows by Corollary~\ref{cor:cdfweak}.
\end{solution}

\begin{solution}{ex:weakexamples}
\idea Weak convergence controls integrals of bounded continuous functions only; $x\mapsto x$ is unbounded, and a vanishing amount of mass far away can carry a non-vanishing amount of mean. A uniform bound on a higher
moment forbids this, exactly as uniform integrability did in Chapter~8.
\writeup (a) The distribution function of $\mathcal N(0,\frac1n)$ at $t$ is $\Phi(t\sqrt n)$, which tends to $0$ for $t<0$ and $1$ for $t>0$, the distribution function of $\delta_0$ at its continuity points. For the uniform laws on
$\{\frac kn\}$, $F_n(t)=\frac{\lfloor nt\rfloor}n\to t$ on $[0,1]$. (b) $\mu_n=(1-\frac1n)\delta_0+\frac1n\delta_n\Rightarrow\delta_0$ (for bounded continuous $f$, $\int f\,d\mu_n=(1-\frac1n)f(0)+\frac1nf(n)\to f(0)$), but $\int x\,d\mu_n=1$ for all
$n$. (c) Let $C=\sup_n\int x^2d\mu_n$ and $\tau_M(x)=\max(-M,\min(x,M))$, bounded and continuous. Since $|x-\tau_M(x)|\le\frac{x^2}M$, $|\int x\,d\mu_n-\int\tau_M\,d\mu_n|\le\frac CM$; also $\int x^2d\mu\le C$ (apply weak convergence to
$\min(x^2,K)$ and let $K\to\infty$), so the same bound holds for $\mu$. Hence $\limsup_n|\int x\,d\mu_n-\int x\,d\mu|\le\frac{2C}M$ for every $M$.
\end{solution}

\section*{Sources and further reading}

\begin{itemize}
\item F.~Hausdorff, ``Dimension und \"au\ss eres Ma\ss'', \emph{Math. Ann.} 79 (1918), 157--179.
\item B.~Mandelbrot, ``How long is the coast of Britain? Statistical self-similarity and fractional dimension'', \emph{Science} 156 (1967), 636--638.
\item F.~Riesz, ``Sur les op\'erations fonctionnelles lin\'eaires'', \emph{C. R. Acad. Sci. Paris} 149 (1909), 974--977.
\item Yu.~V.~Prokhorov, ``Convergence of random processes and limit theorems in probability theory'', \emph{Theory Probab. Appl.} 1 (1956), 157--214.
\item A.~N.~Kolmogorov, \emph{Grundbegriffe der Wahrscheinlichkeitsrechnung}, Springer, 1933.
\item P.~Billingsley, \emph{Convergence of Probability Measures}, 2nd ed., Wiley, 1999, Chapter~1. The portmanteau theorem, tightness and Prokhorov.
\item K.~Falconer, \emph{Fractal Geometry: Mathematical Foundations and Applications}, 3rd ed., Wiley, 2014, Chapters~2--4 and~9. Hausdorff measures, the mass distribution principle, self-similar sets.
\item W.~Rudin, \emph{Real and Complex Analysis}, 3rd ed., McGraw--Hill, 1987, Chapter~2. The Riesz representation theorem, whose proof we followed.
\item \emph{For the stories in the introduction.} The table is in L.~F.~Richardson, ``The problem of contiguity: an appendix to \emph{Statistics of Deadly Quarrels}'', \emph{General Systems Yearbook}
6 (1961), 139--187, is the table; Mandelbrot's paper, cited above, made it famous.
\end{itemize}
